\chapter{Einleitung}
\label{cha:Einleitung}

\section{Motivation und Problemstellung}
Die Landwirtschaft steht vor wachsenden Herausforderungen: Der Klimawandel führt zu erhöhtem Schädlingsdruck, veränderten Niederschlagsmustern und sinkenden Erträgen \autocite{calvinIPCC2023Climate2023, hultgrenImpactsClimateChange2025, ClimateChangeTrends}, während eine wachsende Weltbevölkerung und zunehmender Fachkräftemangel den Produktionsdruck zusätzlich verschärfen \autocite{foodandagricultureorganizationoftheunitednationsFutureFoodAgriculture2017, WorldPopulationProspects, duckettAgriculturalRoboticsFuture2018}. Im Rahmen der digitalen Landwirtschaft bieten Drohnen (unbemannte Luftfahrzeuge) vielversprechende Lösungsansätze \autocite{bassoDigitalAgricultureDesign2020}: Sie ermöglichen unter anderem die gezielte Ausbringung von Pflanzenschutzmitteln sowie großflächiges Crop-Monitoring mittels Remote~Sensing \autocite{mashoriRemoteSensingUAVs2026, radoglou-grammatikisCompilationUAVApplications2020} und können so dazu beitragen, die Landwirtschaft effizienter und resilienter zu gestalten \autocite{undeExpandingRoleMultirotor2025, guebsi2024drones}.

Trotz umfangreicher Forschung zu Drohnensystemen bleibt autonomes Landen eine offene Herausforderung \autocite{amendolaDroneLandingReinforcement2024}. Unfallfreies Landen ist für die meisten Flugmissionen essentiell; gleichzeitig erfordern die vielfältigen Landeszenarien jeweils angepasste Lösungen. Bisherige Ansätze reichen von Global Navigation Satellite System (GNSS)-gestützter Positionierung \autocite{patrikGNSSbasedNavigationSystems2019} über markerbasierte visuelle Verfahren \autocite{wubbenAccurateLandingUnmanned2019} bis hin zu markerfreien Systemen, die Landezonen eigenständig aus Sensordaten ableiten \autocite{bartolomeiAutonomousEmergencyLanding2022}. Dabei zeigt sich eine wiederkehrende Einschränkung: Systeme zur Hindernisvermeidung navigieren erfolgreich durch komplexe Umgebungen, ohne gezielt zu landen \autocite{loquercioLearningHighspeedFlight2021, xuNavRLLearningSafe2024}, während Landesysteme überwiegend in hindernisfreien Szenarien operieren \autocite{polvaraSimtoRealQuadrotorLanding2020, ladoszAutonomousLandingMoving2024}. Die Kombination beider Teilaufgaben, Präzisionslandung mit gleichzeitiger Hindernisvermeidung, ist insbesondere in natürlichen Umgebungen mit dichter Vegetation bislang kaum untersucht. Weinberge sind ein konkretes Beispiel solcher Umgebungen: Reihenabstände von typischerweise 1,5 bis 2\,m und dichtes Blattwerk erzeugen beengte Platzverhältnisse, die eine autonome Landung zwischen den Reihen erheblich erschweren.

Reinforcement Learning (RL) hat sich in jüngerer Zeit als wirksame Lösungsstrategie für Drohnenaufgaben etabliert \autocite{kaufmannChampionlevelDroneRacing2023, hodgeDeepReinforcementLearning2021, amendolaDroneLandingReinforcement2024}: Ein RL-Agent erlernt eine Steuerungsstrategie durch Interaktion mit einer simulierten Umgebung, ohne dass ein explizites Umgebungsmodell oder manuell definierte Steuerungsregeln benötigt werden. In Kombination mit Convolutional Neural Networks (CNNs) können aufgabenrelevante Merkmale direkt aus Sensordaten extrahiert werden \autocite{mnihHumanlevelControlDeep2015}, was eine bildbasierte Hindernisvermeidung ohne manuelle Merkmalsdefinition ermöglicht.

Da Drohnen engen Gewichts- und Energiebeschränkungen unterliegen \autocite{semerikovVisionBasedAutonomousUAV2025}, beschränkt sich die vorliegende Arbeit auf eine minimale Sensorkonfiguration aus Tiefenkamera und propriozeptivem Zustandsvektor und setzt auf einen lernbasierten Ansatz, um diese begrenzte Sensorik für die autonome Landung in hindernisreichen Umgebungen auszunutzen.

\section{Zielstellung}

Ziel der Arbeit ist die Entwicklung eines Reinforcement-Learning-basierten Verfahrens, das eine autonome Multirotor-Drohne befähigt, mithilfe lokaler Tiefensensorik in einer hindernisreichen Umgebung sicher zu landen. Das System kombiniert einen propriozeptiven Zustandsvektor mit tiefenbildbasierter Umgebungswahrnehmung, um Hindernisse während des Landevorgangs zu erkennen und ihnen auszuweichen. Die Leistungsfähigkeit des trainierten Agenten wird anhand einer abstrakten Trainingsumgebung evaluiert und anschließend im Zero-Shot-Transfer, also ohne erneutes Training in der Zielumgebung, auf eine Umgebung mit realistischer Vegetationsstruktur getestet.

Dadurch ergeben sich folgende Fragestellungen:

\minisec{Hauptfragestellung}
Inwiefern kann ein Reinforcement-Learning-basierter Ansatz eine autonome Drohne befähigen, mithilfe lokaler Tiefensensorik in hindernisreichen Umgebungen sicher zu landen?

\minisec{Nebenfragestellungen}
\begin{enumerate}
	\item Wie leistungsfähig ist der trainierte Agent in der Trainingsumgebung, und welche Faktoren begrenzen die Erfolgsrate?
	\item In welchem Maß lässt sich die erlernte Navigationsleistung auf eine geometrisch verschiedenartige Umgebung mit realistischer Vegetationsstruktur übertragen?
	\item Welche Leistungslücke besteht zwischen einem lernbasierten Agenten mit lokaler Sensorik und einem Pfadplaner mit vollständiger Umgebungsinformation?
\end{enumerate}

\section{Aufbau der Arbeit}

Kapitel~\ref{cha:Fachliche Grundlagen} führt die für diese Arbeit relevanten Grundlagen ein: die theoretischen Grundlagen des Reinforcement Learning von der MDP-Formulierung über Wertfunktionen bis hin zu Proximal Policy Optimization, Convolutional Neural Networks zur Merkmalsextraktion aus Bilddaten sowie achsenparallele Bounding Boxes zur Kollisionserkennung. Kapitel~\ref{cha:Stand der Forschung} gibt einen Überblick über bestehende Ansätze zur autonomen Drohnenlandung und Hindernisvermeidung und identifiziert die Forschungslücke, die diese Arbeit adressiert. Kapitel~\ref{cha:Methodik} beschreibt das entwickelte System, einschließlich der MDP-Formulierung, der CNN-MLP-Architektur, der simulierten Trainingsumgebung und des zweiphasigen Evaluationsprotokolls. Kapitel~\ref{cha:Ergebnisse} präsentiert die experimentellen Ergebnisse beider Evaluationsphasen sowie den Vergleich mit einem klassischen Pfadplaner. Kapitel~\ref{cha:Diskussion} interpretiert die Ergebnisse im Kontext der Forschungsfragen und diskutiert Limitationen des Ansatzes. Kapitel~\ref{cha:Fazit und Ausblick} fasst die zentralen Erkenntnisse zusammen und skizziert Perspektiven für weiterführende Arbeiten.
